\subsection{An order-four action on rank-one components}
\begin{proposition}\label{prop:r108}
The ring $R8\_2\_68$ has no complex fusion categorification.
\end{proposition}
\begin{proof}
Its exact type is $(1^4,4^2,6^2)$, its invertibles form $C_4$, its
adjoint is the whole ring, and its order-two invertible fixes all four
noninvertibles. Solvability in dimension $108$ therefore forces
$\mathcal C=\mathcal D^{C_2}$ with $\mathcal D$ integral of dimension
$54$ and type $(1^2,2^4,3^4)$. The outer involution exchanges both
noninvertible dimension classes in pairs and fixes the invertibles.
The possible proper dimensions of its invariant adjoint are
$1,2,9,18,27$. Pointed dimensions one and two cannot contain the square
of a dimension-three object; dimensions nine and twenty-seven would
contain dimension-two objects, contradicting strong Frobenius
divisibility \cite{ENOSolvable}. Thus either
$\mathcal D_{\rm ad}=\mathcal D$ or its adjoint $\mathcal E$ has type
$(1^2,2^4)$ and dimension eighteen.

In the first case prime descent gives
$\mathcal D=\mathcal F^{C_2}$, with $\mathcal F$ of type $(1^9,3^2)$.
In the second case $U(\mathcal E)=1$: a grading of order two would
have a dimension-nine neutral containing a dimension-two simple,
again contradicting strong Frobenius; order-three components of
dimension six would each require both invertibles; higher grading
orders cannot accommodate a dimension-two simple. Hence
$\mathcal E=(\Vec_A^\omega)^{C_2}$, $|A|=9$, with inversion on $A$.
Its unique nontrivial central invertible $g$ is the canonical
Tannakian lift. The relative-center $C_3$ action of the grading of
$\mathcal D$ fixes this lift. Since $H^2(C_3,\C^\times)=0$, its
central equivariant lifts, together with the unit lifts, form an
abelian group of order six. The unique order-two lift restricts to
$g$ and inherits its self-braiding $+1$; it therefore gives a central
$\Rep(C_2)$ in $\mathcal D$. This explicitly verifies bosonicity,
rather than inferring it from centrality alone.
De-equivariantization again gives $\mathcal F$ of type $(1^9,3^2)$:
a dimension-three object fixed by $g$ would descend to dimension
$3/2$, which is not an algebraic integer.

In both cases the order-two central lift in $\mathcal D$ is unique:
its kernel is trivial when $U(\mathcal D)=1$, and has order three
when $U(\mathcal D)=C_3$. The outer involution preserves this
Tannakian subcategory. Combining the two equivariantizations gives
$\mathcal C=\mathcal F^{C_4}$, because the intermediate symmetric
subcategory has dimension four and its four characters form the
census $C_4$. Let $\beta$ generate this coherent action. Its square
acts by inversion on $A=\operatorname{Inv}(\mathcal F)$ and fixes
the two dimension-three objects.

The adjoint of $\mathcal F$ is $\Vec_A^\omega$, with a $C_3$ grading
of ranks $9,1,1$. Rank-one components force $[\omega]=0$ by their
module pentagon. The action of $\beta$ exchanges these two components;
otherwise the original $\mathcal C$ would inherit a nontrivial
$C_3$ grading. For $A=C_9$ no automorphism squares to inversion.
For $A=C_3^2$, its matrix $B$ satisfies $B^2=-I$, so its minimal
polynomial is the irreducible $t^2+1$ over $\mathbb F_3$ and
$\det(B)=1$. Its tensorator class
$\eta\in H^2(A,\C^\times)=C_3$ is fixed by $B$; order-four
coherence gives $4\eta=0$, hence $\eta=0$.
Such a tensor autoequivalence fixes every rank-one left-module class,
since those classes are $H^2(A,\C^\times)$. But the two nonneutral
components have different left-module classes: an equivalence,
tensored with the inverse of one invertible bimodule, would identify
the rank-nine regular module with a rank-one component.
Thus $\beta$ cannot exchange them, a contradiction. The supplement
checks every source coefficient, adjoint dimension partition and
matrix calculation; the central lift and action tower use
\cite{GNN,DGNO,ENOSolvable}.
\end{proof}
